%! TeX program = lualatex
\documentclass[../main.tex]{subfiles}
\begin{document} \section{Differentiation formulas and properties} \label{lesson:integration-formulas}
Derivatives of sophisticated functions are calculated from derivatives of basic functions.
Allow the Leibniz notation \(\frac{d}{dx} \cdots\) to guide you to the next step. 

\medskip\hrule

\[
  \begin{array}{rccll}
    \multicolumn{4}{c}{\text{\textbf{\hlsupp{Differentiation} Formulas}}} & \hspace{4in} \\[2ex]
    \frac{d}{dx} & {\color{supp} cx        } &=& {\color{main} c} 
                 & \hspace{4in} \\[2ex]
    \frac{d}{dx} & {\color{supp} x^{n}   } &=& {\color{main} n x^{n - 1}} \text{ but only for } {\hlwarn{n \ne 0}}
                 & \hspace{4in} \\[2ex]
    \frac{d}{dx} & {\color{supp} \ln(x)    } &=& {\color{main} x^{-1}}
                 & \hspace{4in} \\[2ex]
    \frac{d}{dx} & {\color{supp} e^{x}     } &=& {\color{main} e^{x}}
                 & \hspace{4in} \\[2ex]
    \frac{d}{dx} & {\color{supp} \sin(x)   } &=& {\color{main} \cos(x)}
                 & \hspace{4in} \\[2ex]
    \frac{d}{dx} & {\color{supp} \cos(x)   } &=& {\color{main} -\sin(x)}
                 & \hspace{4in} \\[2ex]
    \frac{d}{dx} & {\color{supp} \tan(x)   } &=& {\color{main} \sec^{2}(x)}
                 & \hspace{4in} \\[2ex]
    \frac{d}{dx} & {\color{supp} \sec(x)   } &=& {\color{main} \tan(x)\sec(x)}
                 & \hspace{4in} \\[2ex]
    \frac{d}{dx} & {\color{supp} \arctan(x)} &=& {\color{main} \frac{1}{x^{2}+1}}
                 & \hspace{4in} \\[2ex]
    \frac{d}{dx} & {\color{supp} \arcsin(x)} &=& {\color{main} \frac{1}{\sqrt{1 - x^{2}}}} 
                 & \hspace{4in} 
  \end{array}
\]

\hrule
\textbf{Differentiation Rules} (for \(+, -, \times, \div\)).
\hfill{} \faExclamationTriangle{} Only apply to differentiable functions. 

\begin{align*}
  \frac{d}{dx} \bigg(f(x) + g(x)\bigg) &= f'(x) + g'(x) &\text{(sum)} \\[2ex]
  \frac{d}{dx} \bigg(f(x) - g(x)\bigg) &= f'(x) - g'(x) &\text{(difference)}\\[2ex]
  \frac{d}{dx} c \, f(x)  &= c f'(x) & \text{(constant multiple)}\\[2ex]
  \frac{d}{dx} f(x) g(x)  &= f'(x)g(x) + f(x)g'(x) &\text{(product)}\\[2ex]
  \frac{d}{dx} \frac{f(x)}{g(x)}  &= \frac{f'(x)g(x) - f(x)g'(x)}{\big[g(x)\big]^{2}} & \text{(quotient)}\\[2ex]
\end{align*}

\textbf{The Chain Rule} (aka the differentiation rule for composite functions).
\[
  \frac{d}{dx} f(g(x)) = f'\big({\color{supp}g(x)}\big)g'(x)
  \quad\text{or in Leibniz notation,}\quad
  \frac{d}{dx} \underbrace{f\big(\overbrace{\color{supp}g(x)}^{{\color{main}u}}\big)}_{y = f({\color{main}u})}
                                  = \frac{dy}{d{\color{main}u}} \cdot \frac{d{\color{supp}u}}{dx}
\]
where, in Leibniz notation, \(y = f({\color{main}u})\) is the \emph{result} of the change of variable \({\color{main}u} \;=\; {\color{supp}g(x)}\).
The green \({\color{main}u}\) is the name of the change of variable. The purple \({\color{supp}u}\) is a placeholder for the inside function \({\color{supp}g(x)}\).

\medskip
\hrule

\textbf{Implicit differentiation} (a special application of the chain rule).

% If we denote \(y = f(x)\), then \(f'(x)\) and \(\frac{dy}{dx}\) are just different notations of the same thing, the derivative of \(f(x)\).  Therefore, the expression \(F(y)\) is a composition of functions \(F(f(x))\). We just do not have a concrete expression for \(f(x)\). 

Applying implicit differentiation always involves two steps:
\begin{enumerate}
  \item Under the assumption that \(y\) is implicitly a function of \(x\), apply \(\frac{d}{dx}\) to \hlmain{both sides of an equation} using the chain rule. 
  \item Isolate for \(\frac{dy}{dx}\).
\end{enumerate}

Typical (but not all) applications of implicit differentiation include:
\begin{enumerate}
  \item Given a curve defined by an equation, find the equation of the tangent line at a point on the curve.  
  \item Given an invertible function \(f(x)\), find the derivative of \(f^{-1}(x)\) by implicitly differentiating the equivalent equation \(x = f(y)\).
  \item Related rates.
  \item Logarithmic differentiation.
\end{enumerate}

\medskip
\hrule

\textbf{Convenient formulas} (some special cases of the chain rule). You decide whether to use the chain rule on the fly or memorize the following formulas.

\begin{align*}
  \frac{d}{dx} b^{x} &= \ln(b) b^{x}, \text{ for } b > 0 \text{ and } b \ne 1 \\[2ex]
  \frac{d}{dx} \bigg( g(x) \bigg)^{n} &= n \bigg( g(x) \bigg)^{n-1} g'(x), \quad\text{ but only for } n \ne 0 \\[2ex]
  \frac{d}{dx} e^{g(x)} &= e^{g(x)} \cdot g'(x) \\[2ex]
  \frac{d}{dx} \ln(g(x)) &= \frac{g'(x)}{g(x)}
\end{align*}

% \medskip
% \hrule
% Typical (but not all) questions involving derivatives.
% \begin{enumerate}
%   \item Find the derivative of \(f(x)\) at a constant \(a\). 
%   \item Find values of \(x\) at which a function or a curve defined by an equation has a horizontal tangent line.
%   \item Find values of \(x\) at which a function or a curve defined by an equation has a tangent line of slope \ldots{}.
% \end{enumerate}

\end{document}
